%!TEX TS-program = lualatex
%!TEX encoding = UTF-8 Unicode

\showpoints

\question[12]
The table shows the presence and absence of seven homologies (1–7) among eight species (A–H). Use these data to draw a phylogenetic tree. Draw a tick mark across the branch of each common ancestor and label it with the homologous character (from the table) that unites each clade.

%\vspace*{\baselineskip}

\begin{longtable}[l]{@{}lrrrrrrr@{}}
\toprule
	& 1	& 2	& 3	& 4	& 5	& 6	& 7	\tabularnewline
\midrule
A	& 0	& 1	& 1	& 0	& 0	& 1	& 0	\tabularnewline
B	& 0	& 1	& 0	& 1	& 1	& 0	& 0	\tabularnewline
C	& 0	& 1	& 0	& 0	& 1	& 0 & 1 \tabularnewline
D	& 1	& 1	& 1	& 0	& 0	& 1 & 0	\tabularnewline
E	& 1	& 1	& 1	& 0	& 0	& 1	& 0	\tabularnewline
F	& 0	& 1	& 0	& 1	& 1	& 0	& 0	\tabularnewline
G	& 0	& 1	& 1	& 0	& 0	& 0	& 0	\tabularnewline
H	& 0	& 1	& 0	& 0	& 1	& 0	& 1	\tabularnewline
\bottomrule
\end{longtable}

\vspace*{2\baselineskip}

\begin{AnswerPage}{0.4\textheight}

\centering

\begin{forest}
  mytree,
    [[,name=ALL
		[,name=GADE
			[,name=ADE
				[,name=DE
					[D]
					[E]
				]
				[A]
			]
			[G]
		]
		[,name=CFB
			[,name=FB
				[F]
				[B]
			]
			[,name=CH
			 [C]
			 [H]
			]
		]
    ]]
    \draw [thick] ([xshift=-2mm, yshift=-6mm] DE.base)--+(0:4mm) node [right] {1};
    \draw [thick] ([xshift=-2mm, yshift=-6mm] ADE.base)--+(0:4mm) node [right] {6};
    \draw [thick] ([xshift=-2mm, yshift=-5mm] GADE.base)--+(0:4mm) node [right] {3};
    \draw [thick] ([xshift=-2mm, yshift=-5mm] FB.base)--+(0:4mm) node [right] {4};
    \draw [thick] ([xshift=-2mm, yshift=-5mm] CFB.base)--+(0:4mm) node [right] {5};
    \draw [thick] ([xshift=-2mm, yshift=-5mm] ALL.base)--+(0:4mm) node [right] {2};
    \draw [thick] ([xshift=-2mm, yshift=-5mm] CH.base)--+(0:4mm) node [right] {7};
\end{forest}


\end{AnswerPage}

%\question[3]
%Species C, D, and E are yellow. All other species are blue. Based on your tree, write below how many times  yellow color evolved. Mark on your tree where yellow color evolved.
%
%\ifprintanswers\textbf{{Two times, once for C and once in the common ancestor of D and E.}}\fi

%\question[3]
%Based on your tree above, identify two sister species and tell why %they are sisters.

%\AnswerBox{0.1\textheight}{Sister species are a pair of species that share a recent common ancestor.}
